\documentclass[10pt,a4paper]{amsart}
\usepackage{amsmath,amssymb,amsthm,graphicx}
\usepackage{verbatim}
\newenvironment{reference}{\comment}{\endcomment}
%\newenvironment{reference}{}{}
\newenvironment{slogan}{\comment}{\endcomment}
\newenvironment{history}{\comment}{\endcomment}

% For commutative diagrams we use Xy-pic
\usepackage[all]{xy}

% We use 2cell for 2-commutative diagrams.
\xyoption{2cell}
% Theorem environments.
%
\theoremstyle{plain}
\newtheorem{theorem}[subsection]{Theorem}
\newtheorem{proposition}[subsection]{Proposition}
\newtheorem{lemma}[subsection]{Lemma}

\theoremstyle{definition}
\newtheorem{definition}[subsection]{Definition}
\newtheorem{example}[subsection]{Example}
\newtheorem{exercise}[subsection]{Exercise}
\newtheorem{situation}[subsection]{Situation}

\theoremstyle{remark}
\newtheorem{remark}[subsection]{Remark}
\newtheorem{remarks}[subsection]{Remarks}

\numberwithin{equation}{subsection}

% Macros
%
\def\lim{\mathop{\mathrm{lim}}\nolimits}
\def\colim{\mathop{\mathrm{colim}}\nolimits}
\def\Spec{\mathop{\mathrm{Spec}}}
\def\Hom{\mathop{\mathrm{Hom}}\nolimits}
\def\Ext{\mathop{\mathrm{Ext}}\nolimits}
\def\SheafHom{\mathop{\mathcal{H}\!\mathit{om}}\nolimits}
\def\SheafExt{\mathop{\mathcal{E}\!\mathit{xt}}\nolimits}
\def\Sch{\mathit{Sch}}
\def\Mor{\mathop{\mathrm{Mor}}\nolimits}
\def\Ob{\mathop{\mathrm{Ob}}\nolimits}
\def\Sh{\mathop{\mathit{Sh}}\nolimits}
\def\NL{\mathop{N\!L}\nolimits}
\def\CH{\mathop{\mathrm{CH}}\nolimits}
\def\proetale{{pro\text{-}\acute{e}tale}}
\def\etale{{\acute{e}tale}}
\def\QCoh{\mathit{QCoh}}
\def\Ker{\mathop{\mathrm{Ker}}}
\def\Im{\mathop{\mathrm{Im}}}
\def\Coker{\mathop{\mathrm{Coker}}}
\def\Coim{\mathop{\mathrm{Coim}}}

% Boxtimes
%
\DeclareMathSymbol{\boxtimes}{\mathbin}{AMSa}{"02}

%
% Macros for moduli stacks/spaces
%
\def\QCohstack{\mathcal{QC}\!\mathit{oh}}
\def\Cohstack{\mathcal{C}\!\mathit{oh}}
\def\Spacesstack{\mathcal{S}\!\mathit{paces}}
\def\Quotfunctor{\mathrm{Quot}}
\def\Hilbfunctor{\mathrm{Hilb}}
\def\Curvesstack{\mathcal{C}\!\mathit{urves}}
\def\Polarizedstack{\mathcal{P}\!\mathit{olarized}}
\def\Complexesstack{\mathcal{C}\!\mathit{omplexes}}
% \Pic is the operator that assigns to X its picard group, usage \Pic(X)
% \Picardstack_{X/B} denotes the Picard stack of X over B
% \Picardfunctor_{X/B} denotes the Picard functor of X over B
\def\Pic{\mathop{\mathrm{Pic}}\nolimits}
\def\Picardstack{\mathcal{P}\!\mathit{ic}}
\def\Picardfunctor{\mathrm{Pic}}
\def\Deformationcategory{\mathcal{D}\!\mathit{ef}}
\usepackage{cleveref}
\crefname{originaldefinition}{original definition}{original definitions}
\Crefname{originaldefinition}{Original definition}{Original definitions}
\makeatletter
\newlabel{definition-local-source-target}{{06F8}{source}{}{sourcepointer-1}{}}
\newlabel{definition-local-source-target@cref}{{[originaldefinition][0][]06F8}{[1][0][]source}{}{}{}}
\makeatother
\crefname{originallemma}{original lemma}{original lemmas}
\Crefname{originallemma}{Original lemma}{Original lemmas}
\makeatletter
\newlabel{lemma-local-source-target-implies}{{06F9}{source}{}{sourcepointer-2}{}}
\newlabel{lemma-local-source-target-implies@cref}{{[originallemma][0][]06F9}{[1][0][]source}{}{}{}}
\makeatother
\title{Seven equivalent characterizations of smooth local on the source-and-target}
\author{The Stacks Project Authors}
\date{stacks-project:spaces-descent.tex:local-snapshot-20261012; GNU FDL 1.2 or later}
\begin{document}
\maketitle
\noindent\textit{Editorial note.} All delivered mathematical text is the author's own, copied verbatim from the pinned Stacks Project chapter file spaces-descent.tex; no statement or proof step was rewritten, reordered or removed.\par\medskip
\noindent\textit{Editorial note.} The counted span cites the same-chapter definition and implies-lemma, recorded in assembly.external\_references with canonical tags as navigation pointers.\par\medskip
\noindent\textit{Allowed standard prerequisites (editorial).}\begin{itemize}
\item Definition of smooth local on the source-and-target. Stacks Project, Descent on Spaces, Definition definition-local-source-target. Supplies the (a)-(b) equivalence.
\item Consequences of smooth locality on the source-and-target. Stacks Project, Descent on Spaces, Lemma lemma-local-source-target-implies. Supplies (a)-(e) and (a)-(f), hence (a)-(g).
\item Smooth base change and precomposition stability. Stacks Project, Descent on Spaces, Lemmas lemma-pullback-property-local-target and lemma-precompose-property-local-source. Supplies the (d) equivalence steps.
\end{itemize}\medskip
\ifdefined\hypertarget\hypertarget{sourcepointer-1}{}\fi
\noindent\textit{Source reference.} Original definition 06F8: sources/stacks-chow/tags.txt, source lines 8242--8242. Cross-chapter standard prerequisite cited by the author inside the counted span; it is used as a pointer and its own proof is not part of this item.\par\medskip
\ifdefined\hypertarget\hypertarget{sourcepointer-2}{}\fi
\noindent\textit{Source reference.} Original lemma 06F9: sources/stacks-chow/tags.txt, source lines 8243--8243. Cross-chapter standard prerequisite cited by the author inside the counted span; it is used as a pointer and its own proof is not part of this item.\par\medskip
\begin{lemma}
\label{lemma-local-source-target-characterize}
Let $S$ be a scheme.
Let $\mathcal{P}$ be a property of morphisms of algebraic spaces over $S$
which is smooth local on source-and-target. Let $f : X \to Y$ be a morphism
of algebraic spaces over $S$. The following are equivalent:
\begin{enumerate}
\item[(a)] $f$ has property $\mathcal{P}$,
\item[(b)] for every $x \in |X|$ there exists a smooth morphism of pairs
$a : (U, u) \to (X, x)$, a smooth morphism $b : V \to Y$, and
a morphism $h : U \to V$ such that $f \circ a = b \circ h$ and
$h$ has $\mathcal{P}$,
\item[(c)] for some commutative diagram
$$
\xymatrix{
U \ar[d]_a \ar[r]_h & V \ar[d]^b \\
X \ar[r]^f & Y
}
$$
with $a$, $b$ smooth and $a$ surjective the morphism $h$ has $\mathcal{P}$,
\item[(d)] for any commutative diagram
$$
\xymatrix{
U \ar[d]_a \ar[r]_h & V \ar[d]^b \\
X \ar[r]^f & Y
}
$$
with $b$ smooth and $U \to X \times_Y V$ smooth
the morphism $h$ has $\mathcal{P}$,
\item[(e)] there exists a smooth covering $\{Y_i \to Y\}_{i \in I}$ such
that each base change $Y_i \times_Y X \to Y_i$ has $\mathcal{P}$,
\item[(f)] there exists a smooth covering $\{X_i \to X\}_{i \in I}$ such
that each composition $X_i \to Y$ has $\mathcal{P}$,
\item[(g)] there exists a smooth covering $\{Y_i \to Y\}_{i \in I}$ and
for each $i \in I$ a smooth covering
$\{X_{ij} \to Y_i \times_Y X\}_{j \in J_i}$ such that each morphism
$X_{ij} \to Y_i$ has $\mathcal{P}$.
\end{enumerate}
\end{lemma}

\begin{proof}
The equivalence of (a) and (b) is part of
Definition \ref{definition-local-source-target}.
The equivalence of (a) and (e) is
Lemma \ref{lemma-local-source-target-implies} part (2).
The equivalence of (a) and (f) is
Lemma \ref{lemma-local-source-target-implies} part (1).
As (a) is now equivalent to (e) and (f) it follows that
(a) equivalent to (g).

\medskip\noindent
It is clear that (c) implies (b). If (b) holds, then for any
$x \in |X|$ we can choose a smooth morphism of pairs
$a_x : (U_x, u_x) \to (X, x)$, a smooth morphism $b_x : V_x \to Y$, and
a morphism $h_x : U_x \to V_x$ such that $f \circ a_x = b_x \circ h_x$ and
$h_x$ has $\mathcal{P}$. Then $h = \coprod h_x : \coprod U_x \to \coprod V_x$
with $a = \coprod a_x$ and $b = \coprod b_x$ is a diagram as in (c).
(Note that $h$ has property $\mathcal{P}$ as $\{V_x \to \coprod V_x\}$
is a smooth covering and $\mathcal{P}$ is smooth local on the target.)
Thus (b) is equivalent to (c).

\medskip\noindent
Now we know that (a), (b), (c), (e), (f), and (g) are equivalent.
Suppose (a) holds. Let $U, V, a, b, h$ be as in (d). Then
$X \times_Y V \to V$ has $\mathcal{P}$ as $\mathcal{P}$ is stable under
smooth base change, whence $U \to V$ has $\mathcal{P}$ as $\mathcal{P}$
is stable under precomposing with smooth morphisms. Conversely, if (d)
holds, then setting $U = X$ and $V = Y$ we see that $f$ has $\mathcal{P}$.
\end{proof}

\end{document}
